\documentclass[a4paper,11pt]{article}

\usepackage[a4paper,margin=2cm]{geometry}
\usepackage[T1]{fontenc}
\usepackage[utf8]{inputenc}
\usepackage[ngerman,english]{babel}
\usepackage{lmodern}
\usepackage{microtype}
\usepackage{xcolor}
\usepackage{parskip}
\usepackage{enumitem}
\usepackage[most]{tcolorbox}
\usepackage{titlesec}
\usepackage[hidelinks]{hyperref}
\usepackage{array}
\usepackage{longtable}
\usepackage[table]{xcolor}

\definecolor{HeaderBlueDark}{RGB}{70,110,170}
\definecolor{RowBlueLight}{RGB}{235,243,255}
\definecolor{RowBlueDark}{RGB}{220,233,250}
\definecolor{SoftGray}{RGB}{90,90,90}

\setlength{\parindent}{0pt}
\setlist[itemize]{leftmargin=1.4em,itemsep=0.35em,topsep=0.35em}
\linespread{1.06}
\renewcommand{\arraystretch}{1.35}
\setlength{\tabcolsep}{5pt}

\titleformat{\section}[block]
{\Large\bfseries\color{HeaderBlueDark}}
{}{0pt}{}
\titlespacing{\section}{0pt}{1.3em}{0.8em}

\titleformat{\subsection}[block]
{\large\bfseries\color{HeaderBlueDark}}
{}{0pt}{}
\titlespacing{\subsection}{0pt}{1.0em}{0.6em}

\newtcolorbox{exbox}{enhanced,breakable,colback=RowBlueLight,colframe=RowBlueDark,boxrule=0.4pt,arc=1.5mm,left=1.2mm,right=1.2mm,top=1mm,bottom=1mm,title={Beispiele \textnormal{(Examples)}},fonttitle=\bfseries,colbacktitle=RowBlueDark,coltitle=black}

\NewDocumentEnvironment{grammarentry}{mm+b}{%
    \begin{tcolorbox}[enhanced,breakable,colback=white,colframe=HeaderBlueDark,boxrule=0.6pt,arc=1.5mm,left=1.5mm,right=1.5mm,top=1mm,bottom=1mm,colbacktitle=HeaderBlueDark,coltitle=white,fonttitle=\bfseries,title={#1\hfill\normalfont\footnotesize #2}]
    #3
    \end{tcolorbox}
}{}

\newcommand{\GermanText}[1]{\textbf{Deutsch: }#1\par}
\newcommand{\EnglishText}[1]{{\small\color{SoftGray}\textbf{English: }#1\par}}
\newcommand{\ExampleItem}[2]{\item \textbf{#1}\\{\small\color{SoftGray}#2}}

\title{Substantivierte Adjektive\\\large \textnormal{(Nominalized Adjectives)}}
\date{}

\begin{document}
\maketitle
\tableofcontents
\newpage

\section{Grundidee \textnormal{(Basic Idea)}}
\begin{grammarentry}{Substantiviertes Adjektiv}{Nominalized adjective}
\GermanText{Ein substantiviertes Adjektiv ist ein Adjektiv, das im Satz die Funktion eines Nomens übernimmt. Es wird deshalb großgeschrieben, behält aber die Deklinationsendungen eines Adjektivs.}
\EnglishText{A nominalized adjective is an adjective that takes over the function of a noun in a sentence. It is therefore capitalized, but it keeps adjective declension endings.}
\GermanText{Beispiel: \textbf{der alte Mann} $\rightarrow$ \textbf{der Alte}. Das Nomen \textbf{Mann} wird weggelassen, weil aus dem Kontext klar ist, wer gemeint ist.}
\EnglishText{Example: \textbf{the old man} $\rightarrow$ \textbf{the old one / the old man}. The noun \textbf{man} is omitted because the context makes clear who is meant.}
\end{grammarentry}

\section{Großschreibung \textnormal{(Capitalization)}}
\begin{grammarentry}{Regel}{Rule}
\GermanText{Substantivierte Adjektive werden immer großgeschrieben.}
\EnglishText{Nominalized adjectives are always capitalized.}
\GermanText{Steht dagegen ein echtes Nomen direkt nach dem Adjektiv, bleibt das Adjektiv kleingeschrieben.}
\EnglishText{If an actual noun directly follows the adjective, the adjective remains lowercase.}
\end{grammarentry}
\begin{exbox}
\begin{itemize}
\ExampleItem{Ich möchte etwas \textbf{Süßes}.}{I would like something sweet.}
\ExampleItem{Ich möchte etwas \textbf{süßes Gebäck}.}{I would like some sweet pastry.}
\ExampleItem{Das \textbf{Wichtigste} ist die Gesundheit.}{The most important thing is health.}
\end{itemize}
\end{exbox}

\section{Deklinationsprinzip \textnormal{(Declension Principle)}}
\begin{grammarentry}{Adjektiv bleibt Adjektiv}{The adjective keeps adjective endings}
\GermanText{Obwohl das Wort wie ein Nomen verwendet wird, wird es weiterhin wie ein Adjektiv dekliniert. Die Endung hängt von Kasus, Genus, Numerus und dem Artikelwort ab.}
\EnglishText{Although the word is used like a noun, it is still declined like an adjective. The ending depends on case, gender, number, and the determiner.}
\GermanText{Darum sagt man zum Beispiel: \textbf{der Alte}, \textbf{den Alten}, \textbf{dem Alten}, \textbf{des Alten}.}
\EnglishText{That is why, for example, one says: \textbf{der Alte}, \textbf{den Alten}, \textbf{dem Alten}, \textbf{des Alten}.}
\end{grammarentry}

\section{Maskulin, Feminin, Neutrum und Plural \textnormal{(Masculine, Feminine, Neuter and Plural)}}
\rowcolors{2}{RowBlueLight}{RowBlueDark}
\begin{longtable}{>{\bfseries}p{2.5cm}|p{2.8cm}|p{2.8cm}|p{2.8cm}|p{2.8cm}}
\rowcolor{HeaderBlueDark}
\color{white}\textbf{Kasus / Case} &
\color{white}\textbf{Maskulin / Masculine} &
\color{white}\textbf{Feminin / Feminine} &
\color{white}\textbf{Neutrum / Neuter} &
\color{white}\textbf{Plural / Plural} \\
Nominativ & der Alte & die Alte & das Gute & die Alten \\
Akkusativ & den Alten & die Alte & das Gute & die Alten \\
Dativ & dem Alten & der Alten & dem Guten & den Alten \\
Genitiv & des Alten & der Alten & des Guten & der Alten \\
\end{longtable}
\GermanText{Diese Tabelle zeigt die Deklination nach dem bestimmten Artikel.}
\EnglishText{This table shows declension after the definite article.}

\section{Mit unbestimmtem Artikel \textnormal{(With the Indefinite Article)}}
\begin{grammarentry}{Starke Endung am Adjektiv}{Strong ending on the adjective}
\GermanText{Nach \textbf{ein/eine} trägt das substantivierte Adjektiv die Endung, die der unbestimmte Artikel nicht zeigt. Deshalb heißt es zum Beispiel \textbf{ein Deutscher}, aber \textbf{der Deutsche}.}
\EnglishText{After \textbf{ein/eine}, the nominalized adjective carries the ending that the indefinite article does not show. Therefore, for example, one says \textbf{ein Deutscher}, but \textbf{der Deutsche}.}
\end{grammarentry}
\begin{exbox}
\begin{itemize}
\ExampleItem{Ein \textbf{Deutscher} wartet draußen.}{A German man is waiting outside.}
\ExampleItem{Ich sehe einen \textbf{Deutschen}.}{I see a German man.}
\ExampleItem{Sie spricht mit einem \textbf{Deutschen}.}{She is speaking with a German man.}
\end{itemize}
\end{exbox}

\section{Substantivierte Adjektive für Personen \textnormal{(Nominalized Adjectives for People)}}
\begin{grammarentry}{Personenbezeichnungen}{Person references}
\GermanText{Sehr häufig ersetzen substantivierte Adjektive ein Nomen wie \textbf{Mann}, \textbf{Frau}, \textbf{Person}, \textbf{Mensch} oder \textbf{Leute}.}
\EnglishText{Very often, nominalized adjectives replace a noun such as \textbf{man}, \textbf{woman}, \textbf{person}, \textbf{human being}, or \textbf{people}.}
\end{grammarentry}
\begin{exbox}
\begin{itemize}
\ExampleItem{der \textbf{Alte}}{the old man / the old one}
\ExampleItem{die \textbf{Kranke}}{the sick woman / the female patient}
\ExampleItem{die \textbf{Reichen}}{the rich people / the rich}
\ExampleItem{ein \textbf{Bekannter}}{a male acquaintance}
\ExampleItem{eine \textbf{Bekannte}}{a female acquaintance}
\end{itemize}
\end{exbox}

\section{Neutrale substantivierte Adjektive für abstrakte Dinge \textnormal{(Neuter Nominalized Adjectives for Abstract Things)}}
\begin{grammarentry}{Das + Adjektiv}{The + adjective}
\GermanText{Mit dem neutralen Artikel \textbf{das} können Adjektive abstrakte Eigenschaften, Ideen oder Dinge bezeichnen.}
\EnglishText{With the neuter article \textbf{das}, adjectives can refer to abstract qualities, ideas, or things.}
\end{grammarentry}
\begin{exbox}
\begin{itemize}
\ExampleItem{das \textbf{Gute}}{the good / what is good}
\ExampleItem{das \textbf{Böse}}{evil / what is evil}
\ExampleItem{das \textbf{Neue}}{what is new / the new thing}
\ExampleItem{das \textbf{Wichtigste}}{the most important thing}
\ExampleItem{das \textbf{Beste}}{the best / the best thing}
\end{itemize}
\end{exbox}

\section{Nach \glqq etwas\grqq{} und \glqq nichts\grqq{} \textnormal{(After ``etwas'' and ``nichts'')}}
\begin{grammarentry}{Wichtiges Muster}{Important pattern}
\GermanText{Nach \textbf{etwas} und \textbf{nichts} steht ein substantiviertes Adjektiv sehr häufig im Neutrum mit der Endung \textbf{-es}.}
\EnglishText{After \textbf{etwas} and \textbf{nichts}, a nominalized adjective very often appears in the neuter with the ending \textbf{-es}.}
\GermanText{Muster: \textbf{etwas + Adjektiv-es} und \textbf{nichts + Adjektiv-es}.}
\EnglishText{Pattern: \textbf{etwas + adjective-es} and \textbf{nichts + adjective-es}.}
\end{grammarentry}
\begin{exbox}
\begin{itemize}
\ExampleItem{etwas \textbf{Neues}}{something new}
\ExampleItem{etwas \textbf{Interessantes}}{something interesting}
\ExampleItem{nichts \textbf{Besonderes}}{nothing special}
\ExampleItem{nichts \textbf{Schlimmes}}{nothing bad}
\end{itemize}
\end{exbox}

\section{Nach \glqq viel\grqq, \glqq wenig\grqq{} und \glqq alles\grqq{} \textnormal{(After ``viel'', ``wenig'' and ``alles'')}}
\begin{grammarentry}{Häufige Strukturen}{Common structures}
\GermanText{Auch nach \textbf{viel}, \textbf{wenig} und \textbf{alles} können substantivierte Adjektive stehen. Die Form hängt von der grammatischen Umgebung ab.}
\EnglishText{Nominalized adjectives can also appear after \textbf{viel}, \textbf{wenig}, and \textbf{alles}. The form depends on the grammatical environment.}
\end{grammarentry}
\begin{exbox}
\begin{itemize}
\ExampleItem{viel \textbf{Gutes}}{a lot of good}
\ExampleItem{wenig \textbf{Neues}}{little that is new}
\ExampleItem{viel \textbf{Interessantes}}{a lot of interesting things}
\ExampleItem{alles \textbf{Gute}}{all the best / everything good}
\ExampleItem{alles \textbf{Notwendige}}{everything necessary}
\end{itemize}
\end{exbox}

\section{Nationalitäten und andere Personenbezeichnungen \textnormal{(Nationalities and Other Person Labels)}}
\begin{grammarentry}{Typisches Deklinationsmuster}{Typical declension pattern}
\GermanText{Bezeichnungen wie \textbf{Deutscher}, \textbf{Deutsche}, \textbf{Arbeitsloser}, \textbf{Erwachsener} oder \textbf{Angestellter} werden wie substantivierte Adjektive dekliniert.}
\EnglishText{Words such as \textbf{Deutscher}, \textbf{Deutsche}, \textbf{Arbeitsloser}, \textbf{Erwachsener}, or \textbf{Angestellter} are declined like nominalized adjectives.}
\end{grammarentry}
\begin{exbox}
\begin{itemize}
\ExampleItem{ein \textbf{Erwachsener}}{an adult man / an adult person}
\ExampleItem{der \textbf{Erwachsene}}{the adult man / the adult person}
\ExampleItem{eine \textbf{Angestellte}}{a female employee}
\ExampleItem{die \textbf{Arbeitslosen}}{the unemployed}
\end{itemize}
\end{exbox}

\section{Substantivierte Partizipien \textnormal{(Nominalized Participles)}}
\begin{grammarentry}{Partizip I und Partizip II}{Present and past participles}
\GermanText{Auch Partizipien können substantiviert werden. Sie werden dann großgeschrieben und wie Adjektive dekliniert.}
\EnglishText{Participles can also be nominalized. They are then capitalized and declined like adjectives.}
\end{grammarentry}
\begin{exbox}
\begin{itemize}
\ExampleItem{der \textbf{Reisende}}{the traveler}
\ExampleItem{die \textbf{Reisenden}}{the travelers}
\ExampleItem{der \textbf{Verletzte}}{the injured man / injured person}
\ExampleItem{die \textbf{Verletzten}}{the injured people}
\ExampleItem{der \textbf{Studierende}}{the student}
\end{itemize}
\end{exbox}

\section{Komparativ und Superlativ \textnormal{(Comparative and Superlative)}}
\begin{grammarentry}{Auch Vergleichsformen können Nomen ersetzen}{Comparative forms can also replace nouns}
\GermanText{Komparative und Superlative können ebenfalls substantiviert werden.}
\EnglishText{Comparatives and superlatives can also be nominalized.}
\end{grammarentry}
\begin{exbox}
\begin{itemize}
\ExampleItem{der \textbf{Ältere}}{the older one}
\ExampleItem{die \textbf{Jüngere}}{the younger one}
\ExampleItem{das \textbf{Beste}}{the best}
\ExampleItem{das \textbf{Schönste}}{the most beautiful thing}
\ExampleItem{das \textbf{Wichtigste}}{the most important thing}
\end{itemize}
\end{exbox}

\section{Wie erkennt man ein substantiviertes Adjektiv? \textnormal{(How to Recognize a Nominalized Adjective)}}
\begin{grammarentry}{Praktischer Test}{Practical test}
\GermanText{Frage dich, ob hinter dem Adjektiv ein Nomen fehlt, das man aus dem Kontext ergänzen kann. Wenn das Adjektiv selbstständig die Rolle dieses Nomens übernimmt, ist es substantiviert.}
\EnglishText{Ask yourself whether a noun is missing after the adjective and can be supplied from the context. If the adjective independently takes over the role of that noun, it is nominalized.}
\end{grammarentry}
\begin{exbox}
\begin{itemize}
\ExampleItem{Ich spreche mit dem \textbf{Neuen}.}{I am speaking with the new one / the new man.}
\ExampleItem{Ich spreche mit dem \textbf{neuen Mitarbeiter}.}{I am speaking with the new employee.}
\end{itemize}
\end{exbox}

\section{Der ursprüngliche Beispielsatz \textnormal{(The Original Example Sentence)}}
\begin{grammarentry}{nichts Besonderes}{nothing special}
\GermanText{Im Satz \glqq Es ist Dienstagmorgen, und hier bedeutet das nichts \textbf{Besonderes}.\grqq{} ist \textbf{Besonderes} ein substantiviertes Adjektiv.}
\EnglishText{In the sentence ``It is Tuesday morning, and here that means nothing \textbf{special},'' \textbf{Besonderes} is a nominalized adjective.}
\GermanText{Es wird großgeschrieben, weil es wie ein Nomen verwendet wird, und es bekommt nach \textbf{nichts} die neutrale Endung \textbf{-es}.}
\EnglishText{It is capitalized because it is used like a noun, and after \textbf{nichts} it takes the neuter ending \textbf{-es}.}
\end{grammarentry}

\section{Typische Fehler \textnormal{(Typical Mistakes)}}
\begin{grammarentry}{Fehler vermeiden}{Avoiding errors}
\GermanText{Falsch: \textbf{nichts besonderes}. Richtig: \textbf{nichts Besonderes}.}
\EnglishText{Wrong: \textbf{nichts besonderes}. Correct: \textbf{nichts Besonderes}.}
\GermanText{Falsch: \textbf{ein Deutsche}. Richtig: \textbf{ein Deutscher}.}
\EnglishText{Wrong: \textbf{ein Deutsche}. Correct: \textbf{ein Deutscher}.}
\GermanText{Falsch: \textbf{der Alt}. Richtig: \textbf{der Alte}.}
\EnglishText{Wrong: \textbf{der Alt}. Correct: \textbf{der Alte}.}
\GermanText{Wichtig: Die Großschreibung ändert nicht die Adjektivdeklination.}
\EnglishText{Important: Capitalization does not change adjective declension.}
\end{grammarentry}

\section{Zusammenfassung \textnormal{(Summary)}}
\begin{grammarentry}{Merksätze}{Key rules}
\GermanText{1. Ein substantiviertes Adjektiv übernimmt die Funktion eines Nomens.}
\EnglishText{1. A nominalized adjective takes over the function of a noun.}
\GermanText{2. Es wird großgeschrieben.}
\EnglishText{2. It is capitalized.}
\GermanText{3. Es wird trotzdem wie ein Adjektiv dekliniert.}
\EnglishText{3. It is nevertheless declined like an adjective.}
\GermanText{4. Nach \textbf{etwas} und \textbf{nichts} steht sehr häufig die neutrale Form auf \textbf{-es}: \textbf{etwas Neues}, \textbf{nichts Besonderes}.}
\EnglishText{4. After \textbf{etwas} and \textbf{nichts}, the neuter form ending in \textbf{-es} is very common: \textbf{etwas Neues}, \textbf{nichts Besonderes}.}
\GermanText{5. Auch Partizipien, Komparative und Superlative können substantiviert werden.}
\EnglishText{5. Participles, comparatives, and superlatives can also be nominalized.}
\end{grammarentry}

\end{document}
